\section{The evaluated one-sided traces and cancellation of the explicit correction}
\label{sec:projected-trace-cancellation}

MA-MO-015 computed the ideal projection without asserting that
its explicit correction vanished.  The actual local-source
return of MA-MO-017 now supplies the additional comparison
needed to evaluate that correction.  The argument below keeps
the source Fourier multiplier and every projection in place.

Retain the nontrivial character and either regime of MA-MO-012.
Use the centered incomplete global multiplier
\[
 A_\eta(z)=E_{\eta,S}(1/2+z)L_\eta(1/2+z),\qquad
 A_{\bar\eta}(z)=E_{\bar\eta,S}(1/2+z)L_{\bar\eta}(1/2+z),
 \tag{MA-MO.360}
\]
and the actual local equivariant-Fourier multiplier
\(m_\eta(z)\) of MA-MO.281.  In particular its finite
conductor, different and Fourier phases are retained.  Put
\[
 \tau_{\eta,S}(g,\phi)=\frac1{2\pi}\int_{\mathbb R}
        \mathcal Mg(it)\frac{m_\eta'(it)}{m_\eta(it)}\,dt
                    +g(0)\int\psi,
       \qquad \psi=1-\phi-\phi^\vee.
 \tag{MA-MO.361}
\]

\paragraph{MA-MO-019 (actual one-sided cancellation and the projected trace).}
For every polynomial exponent \(s\ge0\),
\[
 \begin{aligned}
 \operatorname{Tr}_{J_{\eta,a,s}}B_{\eta,g}
 &=\operatorname{Tr}_{H_{a,s}}T_{\eta,a}(g)\\
 &=\operatorname{Tr}_{H_{-a,s}}T_{\eta,-a}(g)
  =\tau_{\eta,S}(g,\phi).
 \end{aligned}
 \tag{MA-MO.362}
\]
Consequently the correction explicitly defined in (MA-MO.232)
satisfies
\[
                    \mathfrak C_{\eta,a}(g,\phi)=0.
 \tag{MA-MO.363}
\]
This is a consequence of proved source and kernel comparisons.
It is not obtained by omitting the ideal projector.  Indeed,
for the assembled operator \(\mathbb T\) of (MA-MO.188),
\[
 \operatorname{Tr}_{\mathbb H}(\mathbb T\Pi_J)=\tau_{\eta,S},
 \qquad
 \operatorname{Tr}_{\mathbb H}\mathbb T=2\tau_{\eta,S}.
 \tag{MA-MO.364}
\]
The two expressions need not be equal.

\begin{proof}
\emph{The source functional equation with its actual multiplier.}
The nontrivial-sector Poisson equation has zero singular term.
On the original local source it is the identity relating
\(RZ^{\rm geom}_{\bar\eta,S}\mathfrak S_\eta\) to
\(Z^{\rm geom}_{\eta,S}\), as used in MA-MO.181 and
MA-MO.295.  The two actual additive Fourier maps are inverse
on these coinvariants; the global element \(-1\) accounts
for their Fourier-square signs, including the finite factors
in the odd Dirichlet sector.  Since
\(\mathfrak U_\eta=\mathfrak S_{\bar\eta}R\), their
scalar functional equation is exactly
\[
                     A_{\bar\eta}(-z)=A_\eta(z)m_\eta(z).
 \tag{MA-MO.365}
\]
One may obtain it first by Mellin transformation of the actual
Poisson identity on smooth local vectors and then meromorphic
continuation.  Transforms of compact annular source tests
separate each scalar multiplier, so no factor is lost in this
identification.  Equivalently it is the elementary local and
global functional equation with the actual Fourier convention
already fixed in MA-CC/MA-CQ/MA-MO-017.  Replacing the local
map by a bare archimedean transform would not justify (MA-MO.365).

\emph{An exact kernel identity on the one-sided space.}
Fix \(b=a\) or \(-a\).  Write \(z=b+it\), \(w=b+iu\),
and use \(\ell_{b,b}\) from (MA-MO.288).  The full Fourier
kernel \(k_b\) of the assembled operator (MA-MO.187) is
\[
                k_b(t,u)=A_\eta(z)\ell_{b,b}(t,u)
                                             A_\eta(w)^{-1}.
 \tag{MA-MO.366}
\]
Here is a domain-valid verification.  On the compact-frequency
core \(V_b^{-1}C_c^\infty\), the global inverse multiplier
is smooth on the compact input support.  Multiplication by
\(\phi\) preserves weighted Schwartz space, the global
forward multiplier has polynomial growth, and both local
multipliers \(m_\eta^{\pm1}\) have polynomial derivatives
on the input line.  Thus all sequential expressions below
exist on that core.

In the Fourier picture the reflected second source term has
multiplier \(A_{\bar\eta}(-z)\) and the cutoff
\(\phi^\vee\).  Formula (MA-MO.365) therefore gives the
literal algebra
\[
 \begin{aligned}
 T_{\eta,b}(g)
 &=C_g\{1-A_\eta M_\phi A_\eta^{-1}
          -A_{\bar\eta}(-\cdot)M_{\phi^\vee}
                                  A_{\bar\eta}(-\cdot)^{-1}\}\\
 &=A_\eta C_g\{1-M_\phi-m_\eta M_{\phi^\vee}
                         m_\eta^{-1}\}A_\eta^{-1}
   =A_\eta C_g\mathcal L_\eta A_\eta^{-1}.
 \end{aligned}
 \tag{MA-MO.367}
\]
Each Fourier multiplier in this display is taken on the fixed
line.  The local bracket is precisely (MA-MO.283), including
its nonsymmetric residual.  Substitution of its proved kernel
gives (MA-MO.366) on the core.

The right side of (MA-MO.366) is Schwartz: in the gap regime
all reciprocal derivatives are polynomial, while in the
generic regime the Gevrey estimates (MA-MO.294) dominate
the logarithmic-square growth of \(A_\eta(w)^{-1}\).
Multiplication by \(A_\eta(z)\) costs only polynomial
growth.  Hence it defines a bounded trace-class operator.
The compact-frequency core is dense in \(H_{b,0}\), and
both sides agree there.  This proves equality of the actual
bounded operators and of their kernels, without an assertion
that the global inverse is bounded on the whole space.

On the kernel diagonal the actual factors cancel:
\(A_\eta(b+it)A_\eta(b+it)^{-1}=1\).
Equation (MA-MO.297), applied on this line, consequently gives
\[
 \operatorname{Tr}_{H_{b,0}}T_{\eta,b}(g)
   =\frac1{2\pi}\int_{\mathbb R}
      \mathcal Mg(b+it)\frac{m_\eta'(b+it)}{m_\eta(b+it)}\,dt
                              +g(0)\int\psi.
 \tag{MA-MO.368}
\]
The local multiplier is holomorphic and nonzero throughout
\(|\operatorname{Re}z|<1/2\).  Its logarithmic derivative
has polynomial growth on every closed substrip.  The compact
smooth transform \(\mathcal Mg\) decreases faster than
every power there.  Ordinary rectangles therefore move this
integral from \(b\) to zero without residues; their horizontal
integrals tend to zero.  This gives \(\tau_{\eta,S}\).
No global reciprocal is used on the central line.

The same Schwartz kernels give trace-class maps between any
two polynomially weighted versions of \(H_b\).  Composition
with the bounded inclusion from the larger exponent to the
smaller, as in (MA-MO.213)--(MA-MO.214), proves that their
traces are equal.  Thus (MA-MO.368) holds with exponent \(s\).
The actual ideal trace is \(\tau_{\eta,S}\) by the nuclear
source return (MA-MO.284)--(MA-MO.297), already proved on
that completed ideal.  This proves (MA-MO.362).

Now substitute the two evaluated one-sided traces into the
independently proved projector formula (MA-MO.232).  Its first
term is \(\tau_{\eta,S}\), and its left side is also that
value.  Hence (MA-MO.363) follows.  Finally the trace of a
block diagonal operator is the sum of its two block traces,
whereas the projected trace is (MA-MO.190).  This proves
(MA-MO.364).
\end{proof}

\paragraph{A nonzero test, so that the projector still matters.}
Choose \(\delta>0\) smaller than every positive prime logarithm
occurring in the finitely many unramified local factors of
\(m_\eta\); if there are none, choose any fixed positive
\(\delta\).  Let \(g\) be a nonnegative, nonzero compact
Gevrey function supported in \((0,\delta)\).  It satisfies
\(g(0)=0\).  Finite unramified local factors have their
nonconstant Fourier contributions at
\(\pm k\log N\mathfrak p\), and the ramified/conductor
monomials contribute only a multiple of \(g(0)\).  All
such terms therefore vanish on this test.  The real gamma
factor of parity \(e_v\in\{0,1\}\) and the complex gamma
factor give, respectively, the elementary digamma integrals
already derived in MA-CH-001:
\[
 \begin{aligned}
 \tau_{\eta,S}(g,\phi)
  ={}&-\sum_{v\ {
m real}}\int_0^\infty
       \frac{e^{-(e_v+1/2)x}}{1-e^{-2x}}g(x)\,dx\\
     &-\sum_{v\ {
m complex}}\int_0^\infty
       \frac{e^{-x/2}}{1-e^{-x}}g(x)\,dx<0.
 \end{aligned}
 \tag{MA-MO.369}
\]
There is at least one archimedean place.  This proves that
dropping \(\Pi_J\) in (MA-MO.364) would in general double a
nonzero trace, even for the requested test class.  The
example is available unconditionally at every width in the
specified generic/Gevrey set.  It asserts no global
zero-location result.

The operator \(\mathbb T\) in (MA-MO.364) is the assembled
block diagonal operator of MA-MO-012, not a newly constructed
extension of every entry of the original common-space matrix.
The result evaluates the positive-ideal source trace.  The
full narrow zero quotient and the omitted-divisor condition
of MA-FT-002 remain separate.  Its source locators are the
actual Poisson and Fourier equations used in MA-MO-012/017,
the complete source-kernel return (MA-MO.288)--(MA-MO.297),
and the previously derived gamma integrals MA-CH.1 and .6.
All trace equalities use actual operators with compatible
nuclear pairings; no categorical weight or algebraic divisor
class substitutes for that analytic comparison.
